\scanpage{272}
\mspage{P 31}

\deleted{entre le formel et l'algébrique (avec, comme deuxième source d'inspiration,
l'article fondamental GAGA de Serre). Quant à la contribution de Serre à laquelle
je fais allusion dans le texte, il s'agit bien sûr, avant tout, de l'introduction
par lui, en géométrie algébrique abstraite\unsure{,} du point de vue des f\,a\,i\,s\,c\,e\,a\,u\,x
(introduit par J\,e\,a\,n L\,e\,r\,a\,y une douzaine d'années auparavant, dans un
contexte tout différent), dans cet autre article fondamental déjà cité FAC
(``Faisceaux algébriques cohérents'').}

\deleted{A la lumière de ces ``rappels'', si je devais nommer les ``ancêtres'' immédiats
de la nouvelle vision géométrique, ce sont les noms de O\,s\,c\,a\,r Z\,a\,r\,i\,s\,k\,i,
A\,n\,d\,r\,é W\,e\,i\,l, J\,e\,a\,n L\,e\,r\,a\,y et J\,e\,a\,n-P\,i\,e\,r\,r\,e S\,e\,r\,r\,e
qui s'imposent à moi aussitôt. Parmi eux Serre a joué un rôle à part, du fait
que c'est par son intermédiaire surtout que j'ai eu connaissance non seulement
de ses propres idées, mais aussi des idées de Zariski, de Weil et de Leray qui
ont eu à jouer un rôle dans l'éclosion et dans le développement de la géométrie
nouvelle.}

\rule{7em}{0.4pt}

\deleted{Leurs idées, bien sûr, ont joué un rôle important lors de mes premiers pas dans
l'édification de la géométrie arithmétique\quad ; plus, il est vrai, comme points de
départ et comme o\,u\,t\,i\,l\,s (qu'il m'a fallu refaçonner plus ou moins de
toutes pièces, pour les besoins d'un contexte beaucoup plus vaste), que comme
une source d'inspiration qui aurait continué à nourrir\handwritten{\unsure{$ {}^{\sim}$}} mes rêves et mes projets,
au cours des mois et des années. De toutes façons, il était bien clair d'emblée
que, même refaçonnés, ces outils étaient très en deça de ce qui était requis,
pour faire même les tout premiers pas en direction des fantastiques conjectures.}

\deleted{11.\quad Les deux idées-forces cruciales dans le démarrage et dans le développement
de la géométrie nouvelle, ont été celle de s\,c\,h\,é\,m\,a et celle de t\,o\,p\,o\,s.
Apparues à peu près simultanément et en étroite symbiose l'une avec l'autre (*),
elles ont été comme un seul et même n\,e\,r\,f m\,o\,t\,e\,u\,r dans l'essor spectaculaire de la nouvelle géométrie, et ceci dès l'année même de leur apparition.
Pour terminer ce tour d'horizon sur mon oeuvre, il me reste à dire quelque mots
au sujet tout au moins de ces deux idées-là.}

\deleted{La notion de schéma est la plus naturelle, la plus ``évidente'' imaginable,
pour englober en une notion unique la série infinie de notions de ``variété''}

\rule{7em}{0.4pt}

\deleted{(*) Il est question de ce démarrage, qui se place en 1958, dans la note de b. de
p. (*) page 23 . La notion de\unsure{,} site ou de ``t\,o\,p\,o\,l\,o\,g\,i\,e d\,e
G\,r\,o\,t\,h\,e\,n\,d\,i\,e\,c\,k'' (version provisoire de celle de topos) est apparue
dans le sillage immédiat de la notion de schéma. C'est elle à son tour qui fournit le langage nouveau de la ``localisation'' ou de ``la descente'', utilisé à chaque
pas dans le développement du thème et de l'outil schématiques. La notion plus
intrinsèque et plus géométrique de t\,o\,p\,o\,s, restée d'abord implicite au cours
des années suivantes, se dégage surtout à partir de 1963, avec le développement
de la cohomologie étale, et s'impose peu à peu à moi comme la notion la plus fondamentale.}
% archive sheet number: 271
